\documentclass[UTF8]{ctexart}
\usepackage{amsmath}
\usepackage{booktabs}
\usepackage{geometry}
\usepackage{listings}
\usepackage{xcolor}
\usepackage{hyperref}
\usepackage{graphicx}
\usepackage{float}

\geometry{a4paper, left=2.5cm, right=2.5cm, top=2.5cm, bottom=2.5cm}
\hypersetup{colorlinks=true, linkcolor=blue, urlcolor=blue, citecolor=blue}

\lstset{
  basicstyle=\ttfamily\small,
  breaklines=true,
  frame=single,
  keywordstyle=\color{blue},
  commentstyle=\color{gray},
  stringstyle=\color{teal},
  showstringspaces=false,
  tabsize=2,
  columns=flexible,
}

\title{\textbf{系统开发工具基础实验报告（第四周）}}
\author{廖世嘉 \quad 25020007073 \\ 25级计算机科学与技术}
\date{2026年9月14日}

\begin{document}
\maketitle
\tableofcontents
\newpage

\section{实验基本信息}

\begin{table}[htbp]
\centering
\caption{实验基本信息}
\label{tab:info}
\begin{tabular}{ll}
\toprule
项目 & 内容 \\
\midrule
姓名 & 廖世嘉 \\
学号 & 25020007073 \\
专业 & 25级计算机科学与技术 \\
实验环境 & WSL Ubuntu 26.04：课上部分使用 Python/ruff/pytest/Make/curl/jq/LaTeX；\\
 & 课后部分在 C++20/CMake 项目 MudGame 上完成，使用 clang-format、cppcheck、\\
 & pre-commit、lcov、semgrep、actionlint、shellcheck、cargo、latexmk \\
实验日期 & 2026年9月14日 \\
课程仓库 & \href{https://github.com/MrLLLiao/SystemToolkitDevCourse}{https://github.com/MrLLLiao/SystemToolkitDevCourse} \\
练习仓库 & \href{https://github.com/MrLLLiao/MudGame}{https://github.com/MrLLLiao/MudGame}（课后练习所在，仅本地提交，未推送） \\
\bottomrule
\end{tabular}
\end{table}

\section{实验目的}

第四周课上讲代码质量与元编程。课上检查的四道题（q13--q16）按讲义要求，在 greetlab 包上搭本地质量门禁、写 Makefile、用 curl+jq 出报告、最后做一次综合交付。课后练习把同一批工具搬到一个更大的 C++ 项目里跑一遍：给 MudGame 配上格式检查、静态分析、覆盖率、CI、semgrep、自定义 GitHub Action，再补一个 Makefile 和 Rust 版本语义的演示。目的很直接：课上学的工具在真实仓库里是否立得住，哪些地方会卡住，用了才知道。

\section{实验环境与工具}

课上部分沿用前几周的 Python 环境：ruff 0.x、pytest、Make、curl、jq、latexmk。q13--q16 的代码在课程仓库 SystemToolkitDevCourse/week4 下的 q13--q16 目录中。

课后 11 个练习全部在 MudGame 仓库（C++20 / CMake，21 个单元测试）中完成。WSL 侧工具链版本：clang-format 21.1.8、cppcheck 2.19.0、lcov 2.0-1、pre-commit 4.6.2、semgrep 1.176.1、actionlint 1.7.12、shellcheck 0.11.0、cargo/rustc 1.93.1、jq 1.8.1。练习期间发现宿主 git 全局代理被设置成 WSL 内不可达的 127.0.0.1:7890，凡涉及 FetchContent 拉取依赖的命令都临时用 \texttt{GIT\_CONFIG\_COUNT} 覆盖为宿主代理，未改动全局配置。

\section{课上检查内容与结果}

\subsection{实验十三：建立可执行的本地质量门禁}

\subsubsection{实验要求}

复制 q10 为 q13，在 pyproject.toml 中加入 ruff 和 pytest 的最小配置。补充一个正常姓名测试和一个空白姓名测试，每个测试包含有效断言。运行 ruff format、ruff check 和 pytest，修复全部问题，不得全局忽略规则。编写 check.sh，依次执行 ruff format --check、ruff check 和 pytest，运行并保证返回 0。

\subsubsection{配置文件}

pyproject.toml 中新增 ruff 和 pytest 配置：

\begin{lstlisting}
[tool.ruff]
line-length = 88
target-version = "py310"

[tool.ruff.lint]
select = ["E", "F", "W", "I"]

[tool.pytest.ini_options]
pythonpath = ["src"]
testpaths = ["."]
\end{lstlisting}

\subsubsection{测试用例}

test\_cli.py 包含两个测试：

\begin{lstlisting}[language=Python]
import pytest
from greetlab.cli import main

def test_normal_name(capsys):
    import sys
    sys.argv = ["sdt-greet", "--name", "alice"]
    main()
    out = capsys.readouterr().out
    assert "Hello, alice!" in out

def test_blank_name_exits_nonzero():
    import sys
    sys.argv = ["sdt-greet", "--name", "   "]
    with pytest.raises(SystemExit) as exc_info:
        main()
    assert exc_info.value.code == 2
\end{lstlisting}

\subsubsection{质量检查脚本}

check.sh 依次执行三项检查：

\begin{lstlisting}[language=bash]
#!/usr/bin/env bash
set -euo pipefail

echo "=== ruff format --check ==="
ruff format --check .

echo "=== ruff check ==="
ruff check .

echo "=== pytest ==="
python -m pytest test_cli.py -v

echo "=== All checks passed ==="
\end{lstlisting}

\subsubsection{结果}

运行 check.sh 输出：

\begin{lstlisting}
=== ruff format --check ===
4 files already formatted
=== ruff check ===
All checks passed!
=== pytest ===
2 passed
=== All checks passed ===
\end{lstlisting}

退出码为 0，全部检查通过。期间 ruff 发现的导入排序问题已通过 ruff check --fix 修复，未使用全局忽略规则。

\subsection{实验十四：让 Make 只重建真正受影响的产物}

\subsubsection{实验要求}

在 q14 中创建 data.csv、stats.py、build\_report.py 和 report.md 四个源文件，编写 Makefile 包含 all、stats.txt、report.txt 和 clean 目标。准确列出数据文件、脚本和中间产物依赖，将 all 和 clean 声明为 .PHONY。验证首次 make 生成两个产物，第二次无改动时不执行配方，touch data.csv 后重新生成，clean 只删除生成物。

\subsubsection{源文件}

\begin{lstlisting}[language=bash]
# data.csv
name,value
alpha,2
beta,3
gamma,5
\end{lstlisting}

stats.py 读取 data.csv 并计算总和写入 stats.txt：

\begin{lstlisting}[language=Python]
import csv
with open("data.csv") as f:
    total = sum(int(r["value"]) for r in csv.DictReader(f))
open("stats.txt", "w").write(str(total))
\end{lstlisting}

build\_report.py 将 report.md 和 stats.txt 合并为 report.txt：

\begin{lstlisting}[language=Python]
text = open("report.md").read()
total = open("stats.txt").read()
open("report.txt", "w").write(f"{text}\nTotal: {total}\n")
\end{lstlisting}

\subsubsection{Makefile}

\begin{lstlisting}
.PHONY: all clean

all: report.txt

stats.txt: data.csv stats.py
	python3 stats.py

report.txt: report.md stats.txt build_report.py
	python3 build_report.py

clean:
	rm -f stats.txt report.txt
\end{lstlisting}

依赖链：data.csv 和 stats.py 生成 stats.txt，report.md、stats.txt 和 build\_report.py 生成 report.txt。all 依赖 report.txt，clean 删除两个生成物。

\subsubsection{验证结果}

\begin{lstlisting}[language=bash]
# 首次 make
$ make
python3 stats.py
python3 build_report.py

# 第二次 make（无改动）
$ make
make: Nothing to done for 'all'.

# touch data.csv 后
$ touch data.csv && make
python3 stats.py
python3 build_report.py

# clean
$ make clean
rm -f stats.txt report.txt
\end{lstlisting}

首次 make 生成两个产物，第二次 make 因无改动而不执行配方，touch data.csv 后 stats.txt 和 report.txt 均重新生成，clean 只删除生成物，保留源文件。Make 的增量构建机制工作正确。

\subsection{实验十五：把本地 API 数据转换为可读报告}

\subsubsection{实验要求}

在 q15 中创建给定的 packages.json，运行 python -m http.server 8000 启动本地 HTTP 服务。使用 curl -fsS 获取数据，使用 jq 筛选 status 为 active 且 downloads 不少于 100 的记录，按 downloads 降序、name 升序排列。编写 api\_report.sh 将结果生成 summary.md，包含标题和 name/version/downloads 三列表格。

\subsubsection{数据文件}

packages.json 包含 5 条记录：

\begin{lstlisting}
[
  {"name":"alpha","status":"active","downloads":120,"version":"1.2.0"},
  {"name":"beta","status":"inactive","downloads":900,"version":"2.0.0"},
  {"name":"gamma","status":"active","downloads":450,"version":"1.5.1"},
  {"name":"delta","status":"active","downloads":450,"version":"0.9.0"},
  {"name":"epsilon","status":"active","downloads":80,"version":"3.1.0"}
]
\end{lstlisting}

\subsubsection{脚本实现}

api\_report.sh 使用 curl 获取 JSON 数据，通过 jq 进行筛选、排序和 Markdown 表格生成：

\begin{lstlisting}[language=bash]
#!/usr/bin/env bash
set -euo pipefail

URL="http://127.0.0.1:8000/packages.json"
OUTPUT="summary.md"

curl -fsS "$URL" | \
  jq -r '
    ["| name | version | downloads |", "|------|---------|----------|"] as $header |
    [$header[],
     ([.[] | select(.status == "active" and .downloads >= 100)]
      | sort_by(-.downloads, .name)
      | .[]
      | "| \(.name) | \(.version) | \(.downloads) |")
    ] | .[]' > "$OUTPUT"

{
  echo "# Package Summary"
  echo ""
  cat "$OUTPUT"
} > "${OUTPUT}.tmp" && mv "${OUTPUT}.tmp" "$OUTPUT"

echo "Report written to $OUTPUT"
\end{lstlisting}

\subsubsection{结果}

在 q15 目录启动 HTTP 服务后运行 api\_report.sh，生成的 summary.md：

\begin{lstlisting}
# Package Summary

| name | version | downloads |
|------|---------|----------|
| delta | 0.9.0 | 450 |
| gamma | 1.5.1 | 450 |
| alpha | 1.2.0 | 120 |
\end{lstlisting}

筛选结果正确：beta（inactive）和 epsilon（downloads=80 < 100）被排除。delta 和 gamma 的 downloads 均为 450，按 name 升序 delta 排在 gamma 之前。curl -fsS 在服务不可用时返回非零退出码，确保错误可被脚本捕获。

\subsection{实验十六：修复并交付一个陌生的小型工具仓库}

\subsubsection{实验要求}

复制 q13 为 q16 并初始化 Git。把问候语实现临时改成始终返回字面量 Hello, name!，运行 check.sh 确认测试失败，随后定位并修复。编写最小 Makefile 包含 check、build 和 clean，check 调用 check.sh，build 生成 wheel。运行 make check 和 make build，计算 wheel 的 SHA-256，并把最终改动提交为一个内容聚焦的 Git 提交。

\subsubsection{临时破坏与验证}

将 cli.py 中的实现临时改为字面量输出：

\begin{lstlisting}[language=Python]
# 临时改坏
print("Hello, name!")
\end{lstlisting}

运行 check.sh 确认测试失败：

\begin{lstlisting}
=== ruff check ===
F841 Local variable `a` is assigned to but never used
=== pytest ===
FAILED test_normal_name - assert "Hello, alice!" in "Hello, name!"
FAILED test_blank_name_exits_nonzero - Failed: DID NOT RAISE SystemExit
\end{lstlisting}

\subsubsection{修复}

将实现恢复为正确版本：

\begin{lstlisting}[language=Python]
if not a.name.strip():
    raise SystemExit(2)
print(f"Hello, {a.name}!")
\end{lstlisting}

\subsubsection{Makefile}

\begin{lstlisting}
.PHONY: check build clean

check:
	bash check.sh

build:
	python3 -m build --no-isolation

clean:
	rm -rf build dist src/*.egg-info .pytest_cache .ruff_cache
\end{lstlisting}

\subsubsection{构建与交付}

运行 make check 全部通过，运行 make build 成功生成 wheel。计算 SHA-256：

\begin{lstlisting}[language=bash]
$ sha256sum dist/*.whl
f5acb7abe8a6604d88a68cedc447857cf9b7addd89e04c4a154ca489b630d1c6  dist/greetlab_25020007073-0.1.0-py3-none-any.whl
\end{lstlisting}

Git 提交记录：

\begin{lstlisting}[language=bash]
$ git log --oneline
778d461 fix: restore correct greeting and add Makefile with check/build/clean
c4cd828 init: greetlab package from q13
\end{lstlisting}

提交信息包含问题说明、修复内容和 SHA-256 哈希值，内容聚焦于本次修复和构建配置。

\section{课后练习内容与结果}

对应讲义 \href{https://missing-semester-cn.github.io/2026/code-quality/}{code-quality（2026）} 与 \href{https://missing-semester-cn.github.io/2020/metaprogramming/}{metaprogramming（2020）} 两页，共完成 11 个实例。全部在 MudGame 仓库（C++20/CMake）上完成，按练习分步提交，共 18 个 commit，commit 信息格式为 \texttt{feat(xxx): 中文内容}，未推送远端。

\subsection{讲义一：代码质量（code-quality）}

\subsubsection{练习1 配置格式器、linter 与 pre-commit 门禁}

内容：给 C++ 项目配置 clang-format、cppcheck，并接上 pre-commit 钩子，让不合格的代码提交不进去。

做法：新增 \texttt{.clang-format}（BasedOnStyle LLVM，缩进 4 空格，Allman 大括号，对齐仓库现有风格）、\texttt{.cppcheck-suppressions}（unusedFunction、missingIncludeSystem 等文件级抑制）和 \texttt{.pre-commit-config.yaml}（三个 local 钩子：clang-format 检查暂存 C/C++ 文件、cppcheck 检查暂存 src/ 文件、增量构建+ctest）。

结果：全量扫描时发现 1628 处存量格式违规，直接全库格式化会一次产生超大 diff，也不符合"分步提交"的要求，所以采用\textbf{增量采用}策略——门禁只检查 staged 文件，存量问题留给后续逐步清理。钩子第一次真正拦截发生在 time\_service.cpp：cppcheck 报了 performance 建议（passedByValue），加抑制并说明理由后通过。练习期间还发现 clang-format 不能用于 CMake 文件（会把 add\_executable 格式坏），已用 git checkout 恢复。

提交：d789ea9（门禁与钩子）、d8b79c3（time\_service 格式化+抑制）。

\subsubsection{练习2 单元测试与 HTML 覆盖率}

内容：用 lcov + genhtml 生成覆盖率 HTML 报告，找出未覆盖代码并补测试。

做法：独立构建目录 build-cov 加 \texttt{-DCMAKE\_CXX\_FLAGS='--coverage -O0 -g'}，跑完 ctest 后 lcov capture、按 \texttt{*/MudGame/src/*} extract，再 genhtml 输出到 coverage/。gcc 15 的 gcov 数据与 lcov 2.0 存在兼容告警，命令必须带 \texttt{--rc branch\_coverage=1 --ignore-errors mismatch,inconsistent,unused}，否则 capture 直接报错。脚本固化在 \texttt{scripts/coverage.sh}。

结果：初值行覆盖 86.4\%（386/447）、函数 81.1\%（73/90）、分支 52.4\%（297/567）。对照未覆盖行清单，补了两个测试文件：ore\_lookup\_test.cpp 覆盖 Ore/Layer 查询的异常路径（未知矿石/未知矿区抛 out\_of\_range、产出权重回退 0.0）和 requires\_lighting 的真假分支；tool\_upgrade\_test.cpp 覆盖 Tool::upgrade 的满级与升级分支。补充后行覆盖 93.5\%（418/447）、函数 91.1\%（82/90）、分支 60.7\%（344/567）。其中 object.h 的未初始化成员告警被顺带修复（见练习3）。

提交：ff18905。

\subsubsection{练习3 CI 与"故意弄坏"验证}

内容：配置 GitHub Actions 工作流，并故意弄坏代码验证 CI 能抓到。

做法：\texttt{.github/workflows/ci.yml} 两个 job：build-test（安装依赖、cmake 配置、构建、ctest）和 quality（对\textbf{变更文件}跑 clang-format --dry-run --Werror 与 cppcheck，增量策略与本地门禁一致）。用 actionlint 本地校验通过。

结果：本地模拟 CI 的三次破坏：1把 mining\_layers.json 中 shallow 层解锁等级从 1 改成 99 → ctest 失败（ContextGate.StartIdleEntersMining 等用例被数据破坏触发）；2给 object.h 插入错误缩进 → clang-format --dry-run --Werror 返回 1；3还原后 21/21 全绿。另外 cppcheck 全量扫描 src/ 时发现 object.h 的 sellingPrice/buyingPrice 未初始化（默认构造后取值是未定义行为），修复并格式化后单独提交。

提交：1581024（object.h 修复）、6a9e4b3（CI 工作流）。

\subsubsection{练习4 用正则找危险 shell 调用，semgrep 兜底}

内容：按讲义在代码库里搜 \texttt{subprocess.Popen(..., shell=True)} 一类的危险调用，先体验正则在真实代码前的失效方式，再用 semgrep 兜底。

做法与结果（在 /tmp 样例上演示，生产代码本身无 system/popen 调用）：

\begin{lstlisting}[language=bash]
# 正则1（讲义式）：只命中单行调用，多行 Popen 被漏掉
$ grep -nE 'subprocess\.Popen\([^)]*shell\s*=\s*True' demo.py
5:subprocess.Popen("ls -la", shell=True)

# 正则2：os.system 后接行续接符再换行，\bsystem\s*\( 漏报
$ grep -nE '\bsystem\s*\(' demo3.py
（0 命中）

# 破坏正则（去掉 =\s*True）：shell=False 的安全调用被误报
$ grep -nE 'subprocess\.Popen\([^)]*shell' demo.py
16:subprocess.Popen(["ls", "-la"], shell=False)   # 误报

# semgrep 自定义规则：AST 解析，两种危险形态都命中，无误报
$ semgrep --config danger.yml demo.py
Findings: 3 (Popen 单行、Popen 多行、os.system 多行)
\end{lstlisting}

结论：正则写复杂了难维护、写简单了漏报误报；semgrep 按语法树匹配，同样的规则语义下结果稳定。对安全相关检查，工具选型优先级高于正则技巧。

提交：无（技术演示）。

\subsubsection{练习5 正则查找替换项目符号}

内容：把 Markdown 中 ``- '' 项目符号统一替换为 ``* ''。

做法与结果：README.md 先补充了一段用 ``- '' 的功能特性清单，然后用锚定行首的正则批量替换：

\begin{lstlisting}[language=bash]
$ perl -pi -e 's/^\- /\* /' README.md
$ grep -nE '^\- ' README.md   # 验证无残留
$ grep -nE '^\* ' README.md
42:* 采矿系统：5 层矿区（浅层/中层/深层/水晶层/地心），按等级与照明解锁
43:* 时间系统：游戏内时间推进与天气事件联动
44:* 工具系统：3 种工具可升级、修复，耐久随使用消耗
45:* 背包经济：物品交易、背包容量与金币结算
\end{lstlisting}

正则锚定行首、限定单个连字符加空格，不会误伤 \texttt{---} 分隔线或嵌套列表。

提交：1cc8bd9（补充清单）、53b830a（正则替换）。

\subsubsection{练习6 从 JSON 捕获 name 与命令行解析器}

内容：用正则从 JSON 抓 name 字段（处理贪婪量词与转义引号），再写几行代码的命令行解析器。

做法与结果：在含转义引号的样例上对比三种正则（PCRE，grep 的 POSIX ERE 连懒惰量词都不支持）：

\begin{lstlisting}
输入: {"name": "say \"hi\"", "note": "含转义引号"}
贪婪  "name": "(.*)"      -> 捕获到 say \"hi\"", "note": "含转义引号（跨字段，误捕获）
懒惰  "name": "(.*?)"     -> 只捕获到 say \（在转义引号处提前收口，少捕获）
转义感知 "name": "((?:[^"\\]|\\.)*)" -> say \"hi\"（正确，但已难以维护）
\end{lstlisting}

随后用 jq 与几行代码的 python 解析器替代正则：\texttt{scripts/json\_extract.py} 接受 JSON 文件与键名，递归提取所有匹配值。对真实数据 src/Data/Ore/ore.json 提取 name，正确输出黯铁矿、星纹银、赤耀金、雷鸣晶、世界之核 5 个矿石名。

提交：bb2fb43（json\_extract.py）。

\subsection{讲义二：元编程（metaprogramming）}

\subsubsection{练习7 Makefile：git ls-files 驱动与 .PHONY}

内容：按讲义用 git ls-files 驱动文件清单，Makefile 声明 .PHONY 伪目标，避免与同名真实文件冲突。

做法与结果：根目录新增 Makefile，提供 build/test/coverage/list/clean 目标。list 目标用 \texttt{\$(shell git ls-files 'src/*.cpp' 'src/**/*.cpp')} 列出 30 个源文件；build 走 cmake；clean 删除 build/、build-cov/、coverage/。验证 make build、make test（21/21）、make clean 均正常。仓库 .gitignore 原本把 Makefile 当作 CLion 产物忽略，需要加 \texttt{!Makefile} 例外才能提交。

提交：715553e。

\subsubsection{练习8 Rust 版本语义}

内容：用 cargo 实际验证 caret、tilde、通配、范围、精确、多约束六种版本要求的解析结果，并构建一个真实场景（掷骰子）。

做法与结果：rust-version-demo 项目依赖 rand，逐个改写 Cargo.toml 的版本要求并查看解析结果（cargo 1.93.1）：

\begin{table}[htbp]
\centering
\caption{Rust 版本语义实测（rand）}
\label{tab:rust}
\begin{tabular}{lll}
\toprule
版本要求 & 解析结果 & 说明 \\
\midrule
\texttt{\^0.8} & 0.8.8 & caret：$\geq$0.8.0 且 $<$0.9.0（0.x 锁次版本） \\
\texttt{\^0.8.3} & 0.8.8 & caret：$\geq$0.8.3 且 $<$0.9.0 \\
\texttt{\textasciitilde0.8.4} & 0.8.8 & tilde：$\geq$0.8.4 且 $<$0.9.0 \\
\texttt{0.8.*} & 0.8.8 & 通配：仅 0.8.x \\
\texttt{$>$=0.8, $<$0.9} & 0.8.8 & 范围：多约束交集 \\
\texttt{=0.8.4} & 0.8.4 & 精确匹配 \\
\texttt{$>$=1.0, $<$1.5} & 解析失败 & "no matching version" \\
\texttt{$>$=0.8.5} & 0.10.2 & 无上限：直接选最新 \\
\bottomrule
\end{tabular}
\end{table}

构建的骰子场景（rand 0.8 的 gen\_range API）编译通过，连续运行输出点数 6、2、6、5；cargo tree 显示 rand 0.8.8 及其传递依赖。

提交：34f8759（演示项目）、968fb03（提交 Cargo.lock 保证可复现构建）。

\subsubsection{练习9 pre-commit 钩子：构建失败拒绝提交}

内容：让 pre-commit 钩子在构建失败时拒绝提交（对应讲义"make paper.pdf 失败则不提交"）。

做法：练习1 已有 build+ctest 钩子，这次把钩子内部改为走练习7 的 Makefile（make build \&\& make test），让两处复用同一套构建入口。

结果：故意往 object.cpp 追加一行语法错误后提交，pre-commit 依次拦截：cppcheck 报 syntaxError（exit 2），build 钩子编译失败（exit 1），提交被拒绝。还原文件后提交正常通过。演示中踩了一个坑：git restore 只还原工作树不还原暂存区，暂存区仍是坏版本，需要 \texttt{git restore --source=HEAD --staged --worktree} 才能彻底还原。

提交：da08f6f。

\subsubsection{练习10 GitHub Pages 与 shellcheck Action}

内容：给仓库配置 GitHub Pages 部署工作流，并在 CI 中加入 shellcheck 步骤检查 shell 脚本。

做法与结果：shellcheck 0.11.0 本地检查 scripts/ 下 4 个脚本全部通过。ci.yml 的 quality job 增加 shellcheck 步骤（apt 安装后 \texttt{shellcheck -x scripts/*.sh}）。新建 \texttt{.github/workflows/pages.yml}（configure-pages + upload-pages-artifact + deploy-pages）与 docs/index.html 静态站点，推送 master 后自动发布到 Pages（需在仓库 Settings 中把发布源设为 GitHub Actions）。两份工作流都通过 actionlint 校验。

提交：8503c01（shellcheck）、74a7e54（Pages）。

\subsubsection{练习11 自定义 GitHub Action 对 .md 跑 proselint}

内容：写一个自定义复合 Action，对仓库 markdown 文件运行 proselint 并报告文体违规，在 .md 变更时触发。

做法：\texttt{.github/actions/proselint-check/action.yml}（composite action：安装 proselint 后对传入文件运行）；\texttt{md-lint.yml} 计算本次变更的 .md 文件（与 CI 其他检查相同的增量策略）再调用该 Action。

结果：proselint 对 README.md 报告 4 处 ``...'' 应为 ``…'' 的标点问题，修复后本地等价命令通过。期间发现 proselint check 无参数时会读取 stdin 卡死，Action 内加了空文件守卫（无变更则跳过）。

提交：ae2682e（标点修复）、d55253d（Action 与工作流）、10e3774（改为只检查变更文件）。

\subsection{课后练习汇总}

\begin{table}[htbp]
\centering
\caption{11 个课后练习完成情况}
\label{tab:homework}
\begin{tabular}{clll}
\toprule
编号 & 讲义 & 练习主题 & 验证结果 \\
\midrule
1 & code-quality & 格式器+linter+pre-commit 门禁 & 增量门禁生效，两次真实拦截 \\
2 & code-quality & 单元测试与 HTML 覆盖率 & 行覆盖 86.4\%$\to$93.5\%，报告已生成 \\
3 & code-quality & CI 与故意破坏验证 & 数据/格式破坏均被 CI 抓到 \\
4 & code-quality & regex 找危险调用 + semgrep 兜底 & 正则漏报/误报，semgrep 全抓无误报 \\
5 & code-quality & 正则替换项目符号 & 4 行 ``- ''$\to$``* ''，锚定无误伤 \\
6 & code-quality & JSON 捕获 name + 命令行解析器 & 贪婪/懒惰/转义三例对比，jq 与脚本正确 \\
7 & metaprogramming & Makefile（git ls-files + .PHONY） & 30 源文件，build/test/clean 验证通过 \\
8 & metaprogramming & Rust 版本语义 & 8 种约束实测，不可满足正确报错 \\
9 & metaprogramming & pre-commit 拒绝构建失败 & 语法错误提交被双钩子拦截 \\
10 & metaprogramming & GitHub Pages + shellcheck & 4 脚本全过，工作流 actionlint 通过 \\
11 & metaprogramming & 自定义 Action 跑 proselint & 报告 4 处标点问题，修复后通过 \\
\bottomrule
\end{tabular}
\end{table}

\section{解题感悟}

这周是四周里动手最多的一次。课上 q13--q16 的每个工具都小而确定，按讲义一步步走就能过；课后把它们搬到 MudGame 上，才体会到工具链真正的工作量分布：写检查规则本身不难，难在让检查跑通、并且不误伤。

第一个感受是\textbf{门禁要"增量"而不是"全量"}。MudGame 存量代码有 1628 处格式违规，如果按讲义默认做法让 clang-format 检查全库，第一次提交就会被淹没。改成只检查 staged 文件后，每个新提交都是干净的，存量债可以慢慢还。这个取舍在 CI 里同样成立：quality job 只查变更文件，否则 CI 从第一天起就是红的，也就失去意义了。

第二个感受是\textbf{工具链的坑比工具本身更花时间}。几件小事：宿主 git 全局代理指向 WSL 内不可达的 127.0.0.1，FetchContent 拉 CLI11 必失败，最后用 GIT\_CONFIG\_COUNT 临时覆盖；gcc 15 的 gcov 和 lcov 2.0 数据格式不兼容，不加 ignore-errors 覆盖率报告根本出不来；cargo 不认 HTTP\_PROXY 只认 CARGO\_HTTP\_PROXY；proselint 无参数会读 stdin 卡死。每个坑都不深，但排查路径各不相同——读报错、试参数、看文档，一步步来。

第三个感受是\textbf{覆盖率数字要落到"没测到哪"才有意义}。初值行覆盖 86.4\% 看着不低，但分支只有 52.4\%。对照未覆盖行清单发现漏的全是错误路径：未知矿石要抛异常、未知矿区权重要回退 0、工具满级后 upgrade 要返回 false。补上这些用例后分支到 60.7\%，而这些恰好是游戏运行中最容易崩的地方。

第四点是\textbf{分步提交的价值}。18 个 commit 每个只做一件事，门禁日志、覆盖率脚本、CI、semgrep 演示各自独立。中间两次提交被钩子拦下、一次被 cppcheck 语法误报拦下，回看 git log 时每一条都能说清意图，这在一次性全量提交里是做不到的。

最后，练习4 和练习6 让我对正则的态度变了。以前写正则喜欢"一步到位"，现在知道安全相关的匹配要交给按语法树解析的工具（semgrep），数据提取要交给真正的解析器（jq / json 模块），正则适合做定位和快速过滤，不适合做需要精确语义的检查。

\section{问题与解决方法}

\subsection{课上部分}

\subsubsection{ruff 导入排序问题}

在实验十三中，初次运行 ruff check 发现 test\_cli.py 的导入顺序不符合 isort 规则（I001）。使用 ruff check --fix 自动修复后，再次检查全部通过，无需全局忽略规则。

\subsubsection{Makefile 中的 Tab 缩进}

在实验十四中，Makefile 的配方行必须使用 Tab 而非空格缩进。在 Windows 编辑器中保存时，有时会将 Tab 转换为空格，导致 make 报错"missing separator"。解决方法是使用 cat -A 检查 Makefile 中的缩进字符，确保配方行以 \texttt{<TAB>} 开头。

\subsubsection{jq 排序中的多字段排序}

在实验十五中，需要先按 downloads 降序、再按 name 升序排列。jq 的 sort\_by 函数默认按升序排列，对于降序需要在字段前加负号：sort\_by(-.downloads, .name)。这一语法确保 delta 和 gamma 在相同 downloads（450）时按 name 升序排列。

\subsubsection{Git 提交的聚焦性}

在实验十六中，初始提交误将临时文件 commit\_msg.txt 一并提交。通过 git rm --cached 移除该文件并 amend 提交，确保最终提交仅包含 .gitignore 和 Makefile 两个必要文件的变更。

\subsection{课后部分}

\subsubsection{网络代理与依赖拉取}

WSL 内 git 全局代理指向 127.0.0.1:7890（WSL 内不可达），CMake FetchContent 拉取 CLI11/pjh\_json/googletest 全部失败。解决：命令前注入 \texttt{GIT\_CONFIG\_COUNT=1 GIT\_CONFIG\_KEY\_0=http.proxy GIT\_CONFIG\_VALUE\_0=http://172.26.224.1:7890} 覆盖代理（不改全局配置）；cargo 则需设置 CARGO\_HTTP\_PROXY。

\subsubsection{gcov/lcov 版本兼容}

gcc 15 的 gcov 数据与 lcov 2.0 报 mismatch/inconsistent 告警，不带参数时 lcov capture 直接失败。解决：capture/extract 均加 \texttt{--rc branch\_coverage=1 --ignore-errors mismatch,inconsistent,unused}，固化进 scripts/coverage.sh。

\subsubsection{cppcheck 对 gtest 宏的解析缺陷}

cppcheck 2.19 对 gtest 的 \texttt{TEST} 宏内 \texttt{while(成员调用)}、\texttt{EXPECT\_THROW} 等写法报 syntaxError（代码本身编译和运行都正常，已用最小样例复现）。语法错误无法用 suppression 屏蔽，最终把门禁的 cppcheck 限定为只分析生产代码 src/，测试代码由编译器与测试运行保障。

\subsubsection{git restore 的暂存区陷阱}

练习9 演示还原时，git restore 默认只还原工作树，暂存区仍保留坏版本，导致提交继续失败。需要 \texttt{git restore --source=HEAD --staged --worktree} 同时还原两侧。

\subsubsection{proselint 空参数读 stdin}

proselint check 不带文件参数时会从 stdin 读取并阻塞。自定义 Action 与本地验证脚本都必须先判断文件列表为空则跳过。

\section{实验总结}

课上四道题把质量门禁、Make 增量构建、curl+jq 数据处理、综合交付完整走了一遍，流程上没有卡点。课后 11 个练习把它们放到真实的 C++ 仓库里，得到三个结论：第一，质量门禁必须增量生效，否则大仓库第一天就是红的；第二，覆盖率报告的价值在未覆盖清单，而不在百分比本身；第三，正则、静态分析器、构建工具、CI 各有适用边界，选对工具比把单个工具用得更花哨重要。四周的工具在这周汇到了一起：git 分步提交、pre-commit 门禁、CI 自动化、覆盖率与静态分析，从"代码能跑"走到了"改动可验证"。

\section{GitHub 链接与 commit 记录}\label{sec:commits}

课程仓库：\href{https://github.com/MrLLLiao/SystemToolkitDevCourse}{https://github.com/MrLLLiao/SystemToolkitDevCourse}（本报告与 q13--q16 源码所在）。

练习仓库：\href{https://github.com/MrLLLiao/MudGame}{https://github.com/MrLLLiao/MudGame}（11 个课后练习所在，全部改动按练习分步提交，未做单次全量提交）。

MudGame 本次练习的 commit 记录（最新在上）：

\begin{figure}[H]
\centering
\includegraphics[width=\textwidth]{commit_history.png}
\caption{MudGame 课后练习 commit 记录（git log --graph --oneline）}
\label{fig:commits}
\end{figure}

\end{document}
